\documentclass[11pt]{article}

% ---------- packages ----------
\usepackage[utf8]{inputenc}
\usepackage[T1]{fontenc}
\usepackage{lmodern}
\usepackage[a4paper,margin=1in]{geometry}
\usepackage{amsmath,amssymb}
\usepackage{xcolor}
\usepackage[most]{tcolorbox}
\usepackage{booktabs}
\usepackage{array}
\usepackage{enumitem}
\usepackage{titlesec}
\usepackage{parskip}
\usepackage{listings}
\usepackage{graphicx}
\usepackage[colorlinks=true,linkcolor=blue!55!black,urlcolor=blue!60!black,citecolor=blue!55!black]{hyperref}

% ---------- colors ----------
\definecolor{accent}{HTML}{283593}   % indigo
\definecolor{accent2}{HTML}{2E7D32}  % green
\definecolor{accent3}{HTML}{B23A00}  % burnt orange
\definecolor{warn}{HTML}{C62828}     % red
\definecolor{codebg}{HTML}{F5F5F7}
\definecolor{softgray}{HTML}{EEF0F5}

% ---------- section styling ----------
\titleformat{\section}{\Large\bfseries\color{accent}}{\thesection}{0.6em}{}
\titleformat{\subsection}{\large\bfseries\color{accent!85!black}}{\thesubsection}{0.6em}{}
\titleformat{\subsubsection}{\normalsize\bfseries\color{accent!80!black}}{\thesubsubsection}{0.5em}{}
\titlespacing*{\section}{0pt}{1.5em}{0.6em}

% ---------- boxes ----------
\newtcolorbox{primerbox}[1][]{
  colback=accent!6, colframe=accent, boxrule=0.6pt, arc=2pt,
  left=8pt,right=8pt,top=5pt,bottom=5pt, fonttitle=\bfseries, #1}
\newtcolorbox{keybox}[1][]{
  colback=green!5, colframe=accent2, boxrule=0.6pt, arc=2pt,
  left=8pt,right=8pt,top=5pt,bottom=5pt, fonttitle=\bfseries, title={Key idea}, #1}
\newtcolorbox{pitfall}[1][]{
  colback=red!4, colframe=warn, boxrule=0.6pt, arc=2pt,
  left=8pt,right=8pt,top=5pt,bottom=5pt, fonttitle=\bfseries, title={Common pitfall}, #1}
\newtcolorbox{tipbox}[1][]{
  colback=accent3!8, colframe=accent3, boxrule=0.6pt, arc=2pt,
  left=8pt,right=8pt,top=5pt,bottom=5pt, fonttitle=\bfseries, title={Tip}, #1}
\newtcolorbox{goalbox}[1][]{
  colback=softgray, colframe=accent!70!black, boxrule=0.6pt, arc=2pt,
  left=8pt,right=8pt,top=4pt,bottom=4pt, fonttitle=\bfseries, #1}

% ---------- code ----------
\lstdefinestyle{py}{
  language=Python,
  basicstyle=\ttfamily\footnotesize,
  backgroundcolor=\color{codebg},
  keywordstyle=\color{accent}\bfseries,
  commentstyle=\color{gray!60!black}\itshape,
  stringstyle=\color{accent2},
  numbers=left, numberstyle=\tiny\color{gray!70}, numbersep=7pt,
  breaklines=true, breakatwhitespace=false, showstringspaces=false,
  columns=fullflexible, keepspaces=true,
  frame=leftline, framerule=2pt, rulecolor=\color{accent!60},
  xleftmargin=13pt, framexleftmargin=13pt, aboveskip=9pt, belowskip=7pt, tabsize=4,
}
\lstset{style=py}

\newcommand{\term}[1]{\textbf{\color{accent!85!black}#1}}
\newcommand{\R}{\mathbb{R}}
\newcommand{\code}[1]{\texttt{\small #1}}

\begin{document}

% ---------------- TITLE ----------------
\begin{titlepage}
\centering
\vspace*{1.3cm}
{\Huge\bfseries\color{accent} Evaluating Low-Rank / Tensor\\[2pt] Decomposition of the LLaMA-3.2 KV Cache\par}
\vspace{0.5cm}
{\Large A Beginner's Project Handout --- Single Device\par}
\vspace{0.5cm}
{\large Round-trip latency \;\textbullet\; end-to-end latency \;\textbullet\; memory \;\textbullet\; energy \;\textbullet\; accuracy\par}
\vspace{1.3cm}
\begin{tcolorbox}[colback=accent!6,colframe=accent,width=0.93\textwidth,arc=3pt]
\centering
\textbf{What you will be able to say at the end:} ``On LLaMA-3.2-1B/3B, on one GPU, low-rank
(rank-$r$) decomposition of the KV cache compresses keys/values by \underline{X}\%; it costs
\underline{Y} perplexity / \underline{Z}\% task accuracy; and after reconstruction overhead it
is \underline{faster / slower} by \underline{W}\% per token --- here is the compression-vs-accuracy-vs-latency
frontier, and here is which of keys, values, and layers are actually compressible.''
\end{tcolorbox}
\vspace{0.9cm}
\begin{tcolorbox}[colback=softgray,colframe=accent,width=0.93\textwidth,arc=3pt]
\centering\small
Written against your environment: \textbf{transformers 4.57.6, torch 2.8.0}. All architecture
numbers, API calls, and evaluation protocols are verified (sources cited). Code is runnable;
copy it stage by stage.
\end{tcolorbox}
\vfill
{\normalsize Project handout \;\textbullet\; \today\par}
\end{titlepage}

\tableofcontents
\newpage

% ================================================================
\section*{How to use this handout}
\addcontentsline{toc}{section}{How to use this handout}

Read \textbf{Part~I} (concepts) once. Then work \textbf{Part~IV} stage by stage --- each stage
has a goal, the code you need (Part~V), and what to expect (Part~VI). \textbf{Part~III} is your
reference for \emph{how to measure each metric correctly}; skim it before Stage~3 and return to
it whenever you report a number. \textbf{Part~VII} is the literature you are building on. If you
only do one thing first, do \textbf{Stage~1} (the SVD spectrum study) --- it is cheap, needs no
attention surgery, and tells you whether the whole premise even holds on LLaMA-3.2.

% ================================================================
\part*{Part I \;---\; Concepts you need first}
\addcontentsline{toc}{section}{Part I --- Concepts you need first}

\section{The KV cache, exactly}

During generation a transformer produces one token at a time; at each step every layer must
attend to \emph{all previous tokens}. To avoid recomputing them, each layer stores the
\term{keys} (K) and \term{values} (V) of all past tokens --- the \term{KV cache}. It grows by
one entry per layer per generated token, so it grows \emph{linearly with context length} and,
at long context, becomes larger than the model weights (Table~\ref{tab:kv}).

\subsection{Grouped-Query Attention (GQA): get the shape right}
LLaMA-3.2 uses \term{Grouped-Query Attention}: there are many \emph{query} heads but only a few
\emph{key/value} heads, which the query heads share. \textbf{The cache stores the KV heads, not
the query heads} --- this is the \#1 thing beginners get wrong. Per layer, for K and for V:
\[
\text{shape} \;=\; (\text{batch},\; \underbrace{\text{num\_key\_value\_heads}}_{=\,8\text{ for both 1B and 3B}},\; \text{seq\_len},\; \text{head\_dim}).
\]

\begin{table}[h]
\centering
\small
\renewcommand{\arraystretch}{1.2}
\begin{tabular}{@{}lccc@{}}
\toprule
& \textbf{Llama-3.2-1B} & \textbf{Llama-3.2-3B} & (verified from config.json) \\
\midrule
decoder layers & 16 & 28 & \\
hidden size & 2048 & 3072 & \\
query heads & 32 & 24 & \\
\textbf{KV heads (GQA)} & \textbf{8} & \textbf{8} & cache uses this \\
head\_dim & 64 & 128 & \\
FFN size & 8192 & 8192 & \\
vocab & 128256 & 128256 & tied embeddings \\
max context & 131072 & 131072 & \\
default dtype & bfloat16 & bfloat16 & 2 bytes/element \\
\bottomrule
\end{tabular}
\caption{Verified LLaMA-3.2 architecture (base and \texttt{-Instruct} are identical except the EOS token id).}
\label{tab:arch}
\end{table}

\subsection{Per-token KV memory --- the formula you will quote constantly}
\[
\text{KV bytes per token} \;=\; 2 \;(\text{K and V}) \times \text{layers} \times \text{KV heads} \times \text{head\_dim} \times \text{bytes/element}.
\]
For bf16 (2 bytes): \textbf{1B} $= 2\times16\times8\times64\times2 = 32{,}768$ bytes $=32$~KiB/token;
\textbf{3B} $= 2\times28\times8\times128\times2 = 114{,}688$ bytes $=112$~KiB/token.

\begin{table}[h]
\centering
\small
\renewcommand{\arraystretch}{1.2}
\begin{tabular}{@{}lccccc@{}}
\toprule
context (tokens) & 1K & 4K & 8K & 32K & 128K \\
\midrule
1B KV size (32 KiB/tok) & 32 MiB & 128 MiB & 256 MiB & 1.0 GiB & 4.0 GiB \\
3B KV size (112 KiB/tok) & 112 MiB & 448 MiB & 896 MiB & 3.5 GiB & 14.0 GiB \\
\midrule
weights (bf16) & \multicolumn{2}{c}{1B $\approx$ 2.3 GiB} & \multicolumn{3}{c}{3B $\approx$ 6.0 GiB} \\
\bottomrule
\end{tabular}
\caption{KV cache size vs context (batch 1). At 128K the KV cache exceeds the weights; under batched serving it dominates even sooner. \emph{This} is why the KV cache is the thing to compress.}
\label{tab:kv}
\end{table}

\section{What ``tensor / low-rank decomposition'' means here}

\subsection{Matrix low rank (start here)}
Take one layer, one KV head. Its key cache is a matrix $K\in\R^{S\times d}$ ($S$ tokens,
$d=$ head\_dim). A \term{rank-$r$ approximation} writes $K\approx A B^\top$ with $A\in\R^{S\times r}$,
$B\in\R^{d\times r}$ and $r<d$. The best such approximation (in Frobenius norm) comes from the
\term{truncated SVD}: $K=U\Sigma V^\top$, keep the top $r$ singular values.
Storage drops from $Sd$ numbers to $r(S+d)$ --- a win when $r\ll d$. \term{Compression ratio}
$= Sd / (r(S+d))$.

\begin{keybox}
``Low-rank'' works only if the singular values of $K$ decay fast --- i.e.\ a few directions
carry most of the energy. \textbf{Whether that is true for LLaMA's KV is an empirical question}
and is exactly what Stage~1 measures. Spoiler from the literature: keys tend to be low-rank,
\emph{values much less so}~\cite{shadowkv} --- verify it yourself.
\end{keybox}

\subsection{Tensor decompositions, and three granularities}
The full per-layer cache is a 3-D \term{tensor} (heads $\times$ tokens $\times$ head\_dim).
\term{Tucker}, \term{CP/PARAFAC}, and \term{Tensor-Train (TT)} decompositions~\cite{kolda,tt}
compress \emph{several axes jointly} rather than one matrix at a time; the Python library
\code{tensorly} implements all three. You will meet three natural granularities:
\begin{enumerate}[leftmargin=1.6em,itemsep=1pt]
\item \term{per-(layer, head)} --- SVD of each $[S\times d]$ matrix (Stages 1--4). Finest and most interpretable; where the low-rank structure demonstrably is.
\item \term{per-layer, heads flattened} --- SVD of one $[S\times (H\cdot d)]$ matrix per layer; captures cross-head redundancy with a single knob.
\item \term{per-layer 3-D tensor} --- Tucker/TT of $[H\times S\times d]$ (Stage~5); compresses the head mode jointly --- the genuinely ``tensor'' contribution (cf.\ JoLT~\cite{jolt}, DecoQuant~\cite{decoquant}).
\end{enumerate}
\textbf{Start with granularity (1)}, then step up to (3) once it works.

\begin{keybox}
Do all of this \emph{per layer}. The layer axis is for \textbf{rank allocation, not decomposition}:
stack per-layer results only to \emph{budget} ranks across layers (protect the least-compressible
ones), never as a joint decomposition over the layer axis. Different layers are produced by
different projection weights $W_K^{(\ell)}$ and share no common basis, so a cross-layer
factorization saves nothing at attention time (which is per-layer anyway) and only confounds the
depth-dependence you are trying to measure.
\end{keybox}

\section{The one conceptual trap: measure the right error}

The obvious thing to minimize is the \term{reconstruction error} $\lVert K-\hat K\rVert_F$.
\textbf{This is the wrong target.} The cache exists only to produce the attention output
$o=\mathrm{softmax}\!\big(qK^\top/\sqrt d\big)\,V$; a large $\lVert K-\hat K\rVert$ in a direction
the query never looks at is harmless, and a tiny error on a high-attention token can flip the
output. The peer-reviewed statement of this is KDEformer~\cite{kdeformer}: prior methods
minimized $\lVert A-\tilde A\rVert_F$ ``even though $A$ is not the ultimate output of attention
and the output depends on $V$.''

\begin{keybox}
Report \textbf{attention-output error} and \textbf{end-task accuracy}, not just
$\lVert K-\hat K\rVert_F$. Use reconstruction error / the singular-value spectrum as a
\emph{diagnostic} (Stage~1), but judge quality by the output (Stages~2, 4). Recent theory ties
the output error to the \emph{missed attention mass}~\cite{tvattn,adakv}, not to raw KV error.
\end{keybox}

\begin{pitfall}
You may read that ``softmax \emph{amplifies} error when attention is concentrated.'' Taken
literally that is \emph{backwards}: the softmax Jacobian $\mathrm{diag}(p)-pp^\top$ shrinks as
attention peaks, so softmax is \emph{least} sensitive to logit noise when concentrated. The real
reasons output error $\neq$ reconstruction error are (i) value error is weighted by attention
probability ($o=\sum_i p_i v_i$, so $\partial o/\partial v_i = p_i$ --- error on high-attention
tokens matters more); (ii) dropping/mis-estimating attention mass then renormalizing shifts the
output~\cite{tvattn}; and (iii) errors compound over autoregressive steps~\cite{gear}. State it
this way and you will be correct.
\end{pitfall}

% ================================================================
\part*{Part II \;---\; The experiment}
\addcontentsline{toc}{section}{Part II --- The experiment}

\section{Question, scope, variables}

\begin{goalbox}
\textbf{Research question.} On a single device, how does rank-$r$ (and later Tucker/TT)
decomposition of the LLaMA-3.2 KV cache trade \emph{compression} against \emph{accuracy} and
\emph{latency/memory/energy} --- and \emph{which} parts of the cache (keys vs values, shallow vs
deep layers) are actually compressible?
\end{goalbox}

\begin{itemize}[leftmargin=1.4em,itemsep=1pt]
\item \term{Model:} start with \textbf{Llama-3.2-1B} (16 layers --- fast to iterate, fits any
modern GPU), then repeat on \textbf{Llama-3.2-3B} (28 layers --- matches the FSU
poster~\cite{poster}). Use the \code{-Instruct} variants for task accuracy.
\item \term{Baseline:} the unmodified model with a full bf16 KV cache.
\item \term{Method under test:} truncated-SVD rank-$r$ reconstruction of K and V, per (layer, head);
sweep $r$. Extension: Tucker/TT over the (head $\times$ token $\times$ feature) tensor.
\item \term{Independent variables:} rank $r$ (hence compression ratio); context length;
keys-only vs values-only vs both; per-layer rank.
\item \term{The three sub-questions:} \textbf{Q1} does the low-rank premise hold (Stage~1)?
\textbf{Q2} what does it cost accuracy (Stages~2, 4)? \textbf{Q3} is it actually faster/smaller
after overhead (Stage~3)?
\end{itemize}

\begin{tipbox}
Two experiments answer two different questions and must not be conflated. \textbf{Accuracy}
experiments \emph{reconstruct} $\hat K,\hat V$ (dense) and run normal attention --- this measures
quality but saves no memory and \emph{adds} SVD time. \textbf{Memory/latency} experiments must
\emph{store the factors} $A,B$ (memory saved) and reconstruct on the fly (compute added), or
compute attention directly from factors (the research frontier). Keep the two separate; \S\ref{sec:expect}
explains why the naive version can be \emph{slower}.
\end{tipbox}

% ================================================================
\part*{Part III \;---\; Metrics and how to measure them correctly}
\addcontentsline{toc}{section}{Part III --- Metrics and how to measure them}

\section{Definitions and measurement recipes}

\begin{table}[h]
\centering\small
\renewcommand{\arraystretch}{1.25}
\begin{tabular}{@{}p{3.1cm}p{5.4cm}p{5.0cm}@{}}
\toprule
\textbf{Metric} & \textbf{Definition} & \textbf{Tool / pitfall} \\
\midrule
TTFT (time-to-first-token) & latency of the prefill pass (prompt $\to$ first token) & generate 1 token; \code{cuda.synchronize} \\
TPOT / per-step (``round time'') & steady-state cost of one decode step $=(\text{E2E}-\text{TTFT})/(N{-}1)$ & CUDA events; report median \& P95 \\
End-to-end latency & full generation time for $N$ tokens & warmup first; fix seed \\
Throughput & tokens / second & $=N/\text{E2E}$ \\
Peak memory & max GPU bytes during run & \code{max\_memory\_allocated}; reset first \\
KV bytes & cache size & analytic formula \emph{and} measured delta \\
Energy / power & Joules or Watts per token & NVML counter (desktop); \code{tegrastats} (Jetson) \\
Compression ratio & dense bytes / compressed bytes & state keys-only vs both \\
Reconstruction err.\ (diagnostic) & $\lVert K-\hat K\rVert_F/\lVert K\rVert_F$ & \emph{diagnostic only} \\
\textbf{Attention-output err.} & $\lVert o-\hat o\rVert$, or KL of next-token dist. & the error that matters \\
Perplexity & $\exp(\text{mean NLL})$ on WikiText-2 & strided window~\cite{hfppl} \\
Task accuracy & LongBench / RULER / GSM8K & GSM8K guards reasoning \\
\bottomrule
\end{tabular}
\caption{The metric kit. Report accuracy vs compression as a Pareto curve, with latency/memory/energy alongside.}
\end{table}

\begin{pitfall}
\textbf{CUDA is asynchronous.} A bare \code{time.perf\_counter()} around \code{model(...)} times
only the \emph{kernel launch}, not the work --- you will report numbers that are wildly too
small. Always \code{torch.cuda.synchronize()} around wall-clock timing, or use
\code{torch.cuda.Event} (below). Also: do \textbf{warmup} iterations (first calls pay
CUDA-context / cuBLAS / autotune costs), and report \textbf{median and P95}, not mean --- decode
latency is right-skewed.
\end{pitfall}

The measurement helpers (use these everywhere):

\begin{lstlisting}
import torch, numpy as np, time

def time_ms(fn, warmup=10, iters=50):
    """Median/P95 latency of fn() in milliseconds, done correctly."""
    for _ in range(warmup):
        fn()
    torch.cuda.synchronize()
    ts = []
    for _ in range(iters):
        s = torch.cuda.Event(enable_timing=True); e = torch.cuda.Event(enable_timing=True)
        s.record(); fn(); e.record()
        torch.cuda.synchronize()          # REQUIRED before reading elapsed_time
        ts.append(s.elapsed_time(e))
    t = np.array(ts)
    return {"median": float(np.median(t)), "p95": float(np.percentile(t, 95))}

def kv_bytes_measured(model, inputs):
    """Isolate the KV-cache cost of prefill from weights."""
    torch.cuda.synchronize(); torch.cuda.empty_cache()
    torch.cuda.reset_peak_memory_stats()
    base = torch.cuda.memory_allocated()
    with torch.no_grad():
        model(**inputs, use_cache=True)
    torch.cuda.synchronize()
    return torch.cuda.memory_allocated() - base   # includes prefill activations; cross-check w/ formula

def gpu_energy_joules(fn):
    """Desktop NVIDIA energy via the NVML counter (Volta+). pip install nvidia-ml-py"""
    import pynvml
    pynvml.nvmlInit(); h = pynvml.nvmlDeviceGetHandleByIndex(0)
    mj0 = pynvml.nvmlDeviceGetTotalEnergyConsumption(h)  # millijoules
    fn(); torch.cuda.synchronize()
    return (pynvml.nvmlDeviceGetTotalEnergyConsumption(h) - mj0) / 1000.0
\end{lstlisting}

\begin{tipbox}
On a \textbf{Jetson Orin}, energy comes from real on-board sensors: \code{sudo tegrastats --interval 100}
(watch the \code{VDD\_CPU\_GPU\_CV} rail, in mW), or \code{pip install jetson-stats} then read
\code{jtop().power['tot']}. On a desktop GPU use the NVML counter above (more reliable than
\code{nvidia-smi --query-gpu=power.draw} for short kernels).
\end{tipbox}

% ================================================================
\part*{Part IV \;---\; The staged plan}
\addcontentsline{toc}{section}{Part IV --- The staged plan}

\section{Milestones (do them in order)}

\subsection*{Stage 0 --- Setup \& baseline}
\begin{goalbox}
\textbf{Goal:} run the unmodified model and record every baseline metric. \textbf{Produce:} a
table of TTFT, TPOT, E2E, throughput, peak memory, measured KV bytes (vs the analytic formula),
energy/token, and baseline WikiText-2 perplexity, for 1B at a few context lengths.
\end{goalbox}

\subsection*{Stage 1 --- Does the low-rank premise hold? (the cheap, essential first result)}
\begin{goalbox}
\textbf{Goal:} extract K, V per layer from one prefill pass and plot the \emph{singular-value
spectrum} and \emph{energy-retained-vs-rank}. No attention surgery needed. \textbf{Produce:}
(i) spectrum curves per layer; (ii) energy-vs-rank for \emph{keys vs values}; (iii)
energy-vs-rank by \emph{layer depth}. \textbf{Expected (from literature, but verify):} keys much
more compressible than values~\cite{shadowkv}; depth- and RoPE-dependence are \emph{open and
contested}~\cite{shadowkv,lorc} --- a real chance to contribute a clean measurement.
\end{goalbox}

\subsection*{Stage 2 --- Accuracy impact (offline, no attention surgery)}
\begin{goalbox}
\textbf{Goal:} prefill the same prompt with a full cache and with a rank-$r$ cache; compare the
next-token distributions (KL divergence and top-1 agreement) as a function of $r$.
\textbf{Produce:} a KL-vs-$r$ curve; the smallest $r$ preserving the argmax token.
\end{goalbox}

\subsection*{Stage 3 --- Online integration \& real cost}
\begin{goalbox}
\textbf{Goal:} generate with a \code{LowRankCache} and measure TPOT, E2E, throughput, peak
memory, energy. \textbf{Produce:} the honest latency/memory numbers --- \emph{including} whether
naive reconstruction makes decoding \emph{slower} (\S\ref{sec:expect}).
\end{goalbox}

\subsection*{Stage 4 --- Sweep \& frontier}
\begin{goalbox}
\textbf{Goal:} sweep $r$ and context length; add a downstream task (GSM8K for reasoning,
LongBench/RULER for long context). \textbf{Produce:} the compression-vs-accuracy-vs-latency
Pareto frontier; separate keys-only, values-only, both.
\end{goalbox}

\subsection*{Stage 5 --- Baselines \& tensor extension}
\begin{goalbox}
\textbf{Goal:} compare against KIVI 2-bit quantization~\cite{kivi} and the full cache; then try
Tucker/TT (\code{tensorly}) over the per-layer head$\times$token$\times$feature tensor (granularity~3).
\textbf{Produce:} your method placed on the same axes as an established baseline; a first
``genuinely tensor'' result.
\end{goalbox}

% ================================================================
\part*{Part V \;---\; Implementation}
\addcontentsline{toc}{section}{Part V --- Implementation}

\section{Environment}
\begin{lstlisting}[language=bash, keywordstyle=\color{accent}\bfseries]
# Python 3.10+, an NVIDIA GPU. Your versions: transformers 4.57.6, torch 2.8.0
pip install "transformers==4.57.6" torch --upgrade
pip install datasets matplotlib numpy nvidia-ml-py tensorly
# gated model: request access on huggingface.co/meta-llama, then `huggingface-cli login`
# (or use an ungated mirror such as unsloth/Llama-3.2-1B)
\end{lstlisting}

\begin{pitfall}
The KV-cache internal layout is \textbf{version-dependent}. On your 4.57.x, read a layer with the
portable \code{cache[i]} (returns a \code{(K,V)} tuple) and \code{len(cache)} for the layer count;
write in place with \code{cache.layers[i].keys}. Older ($\le$4.53) code used
\code{cache.key\_cache[i]}; the v5 line removes \code{to/from\_legacy\_cache}. Stick to
\code{cache[i]} / \code{len(cache)} in teaching code.
\end{pitfall}

\section{Load the model}
\begin{lstlisting}
import torch
from transformers import AutoTokenizer, AutoModelForCausalLM
NAME = "meta-llama/Llama-3.2-1B-Instruct"     # or unsloth/Llama-3.2-1B-Instruct
tok = AutoTokenizer.from_pretrained(NAME)
model = AutoModelForCausalLM.from_pretrained(NAME, dtype=torch.bfloat16, device_map="cuda").eval()
inputs = tok("Explain why the sky is blue.", return_tensors="pt").to(model.device)
\end{lstlisting}

\section{Stage 1 --- extract K/V and plot the spectrum}
\begin{lstlisting}
import matplotlib.pyplot as plt

with torch.no_grad():
    out = model(**inputs, use_cache=True)      # one prefill pass
cache = out.past_key_values                     # a DynamicCache

def energy_vs_rank(mat):                         # mat: [S, d]
    s = torch.linalg.svdvals(mat.float())        # singular values, descending
    cum = torch.cumsum(s**2, 0) / (s**2).sum()   # fraction of energy kept by top-k
    return cum                                   # cum[r-1] = energy kept at rank r

for i in [0, len(cache)//2, len(cache)-1]:       # shallow / middle / deep layer
    K, V = cache[i]                              # each [B, num_kv_heads, S, head_dim]
    Kmat = K[0, 0]                               # layer i, head 0 -> [S, head_dim]
    Vmat = V[0, 0]
    plt.plot(energy_vs_rank(Kmat).cpu(), label=f"layer {i} K")
    plt.plot(energy_vs_rank(Vmat).cpu(), '--', label=f"layer {i} V")
plt.axhline(0.95, color="gray", ls=":"); plt.xlabel("rank r"); plt.ylabel("energy retained")
plt.legend(); plt.title("Low-rank structure of LLaMA-3.2 KV (keys solid, values dashed)")
plt.savefig("stage1_spectrum.png", dpi=150)
\end{lstlisting}
\emph{Note (RoPE):} the cache stores keys \emph{after} rotary position embedding. So you are
measuring \emph{post-RoPE} low-rankness; the literature disagrees on whether RoPE raises or lowers
key rank~\cite{shadowkv}, so state which you measured. Pre-RoPE analysis is a worthwhile extension.

\section{Stages 2--3 --- a low-rank cache}
The clean, correct pattern: subclass \code{DynamicCache}, and \emph{after the prefill step}
replace the stored keys/values with their rank-$r$ reconstruction; decode tokens append exactly
(a small exact ``tail'').

\begin{lstlisting}
from transformers import DynamicCache

def low_rank_reconstruct(x, r):
    """x: [B, H, S, D] -> rank-r reconstruction of each [S, D] slice (dense, same shape)."""
    B, H, S, D = x.shape
    r = min(r, S, D)
    U, s, Vh = torch.linalg.svd(x.float().reshape(B * H, S, D), full_matrices=False)
    xhat = (U[..., :r] * s[..., :r].unsqueeze(-2)) @ Vh[..., :r, :]
    return xhat.reshape(B, H, S, D).to(x.dtype)

class LowRankCache(DynamicCache):
    def __init__(self, rank=32, compress_keys=True, compress_values=True, **kw):
        super().__init__(**kw)
        self.rank = rank; self.ck = compress_keys; self.cv = compress_values
        self._done = set()
    def update(self, key_states, value_states, layer_idx, cache_kwargs=None):
        keys, values = super().update(key_states, value_states, layer_idx, cache_kwargs)
        if layer_idx not in self._done and keys.shape[-2] > 1:   # once, right after prefill
            if self.ck: keys   = low_rank_reconstruct(keys,   self.rank)
            if self.cv: values = low_rank_reconstruct(values, self.rank)
            self.layers[layer_idx].keys   = keys                 # write back (4.54+)
            self.layers[layer_idx].values = values
            self._done.add(layer_idx)
        return keys, values

# Stage 2: accuracy of the FIRST generated token, full vs rank-r
def next_logits(cache_ctor):
    c = cache_ctor()
    with torch.no_grad():
        o = model(**inputs, past_key_values=c, use_cache=True)
    return o.logits[0, -1].float()

full = next_logits(lambda: DynamicCache(config=model.config))
for r in [4, 8, 16, 32, 64]:
    lr = next_logits(lambda: LowRankCache(rank=r, config=model.config))
    kl = torch.nn.functional.kl_div(lr.log_softmax(-1), full.softmax(-1), reduction="sum")
    agree = (full.argmax() == lr.argmax()).item()
    print(f"r={r:3d}  KL={kl:.4f}  top1_agree={agree}")

# Stage 3: real generation cost with the compressed cache
def gen(cache_ctor, n=64):
    return model.generate(**inputs, do_sample=False, max_new_tokens=n,
                          past_key_values=cache_ctor())
print("dense  ", time_ms(lambda: gen(lambda: DynamicCache(config=model.config))))
print("rank-16", time_ms(lambda: gen(lambda: LowRankCache(rank=16, config=model.config))))
\end{lstlisting}

\begin{tipbox}
\textbf{Memory vs latency, made explicit.} \code{LowRankCache} above reconstructs \emph{dense}
tensors, so it measures \emph{accuracy} and the \emph{reconstruction time} but saves no memory. To
actually save memory, store the factors $A,B$ (and $C,D$ for V) instead of $\hat K$, and
reconstruct inside \code{update()} on read --- memory drops to $r(S+d)$ per head, at the cost of a
reconstruction each step. To save memory \emph{and} avoid reconstruction, compute attention
directly from the factors ($qK^\top\!\approx (qB)A^\top$) --- that is the research frontier
(compressed-domain attention), out of scope here but the natural next step.
\end{tipbox}

\section{Perplexity (baseline quality, Stage 0)}
The canonical strided-window protocol~\cite{hfppl}:
\begin{lstlisting}
from datasets import load_dataset
def perplexity(model, tok, max_length=4096, stride=1024, device="cuda"):
    data = load_dataset("wikitext", "wikitext-2-raw-v1", split="test")
    enc = tok("\n\n".join(data["text"]), return_tensors="pt")
    seq_len = enc.input_ids.size(1); nlls = []; prev = 0
    for begin in range(0, seq_len, stride):
        end = min(begin + max_length, seq_len)
        trg = end - prev
        ids = enc.input_ids[:, begin:end].to(device)
        tgt = ids.clone(); tgt[:, :-trg] = -100          # only last `trg` tokens count
        with torch.no_grad():
            nlls.append(model(ids, labels=tgt).loss)
        prev = end
        if end == seq_len: break
    return torch.exp(torch.stack(nlls).mean()).item()
\end{lstlisting}
This scores the \emph{baseline} model (teacher-forced, no growing cache). Measuring compression's
effect on perplexity requires streaming autoregressive scoring with the compressed cache --- a
worthwhile Stage~4+ extension; for a first pass, judge compression by Stage~2 KL and Stage~4 task
accuracy.

\section{Stage 5 --- per-layer Tucker (granularity 3, the tensor step-up)}
Once the matrix case works, compress the per-layer 3-D tensor $[H\times S\times d]$ with Tucker,
which shares a subspace \emph{across the 8 heads} --- a gain a per-head SVD cannot capture. Use
\code{tensorly} (\code{pip install tensorly}), and keep it \emph{per layer} (never stack the layer
axis --- see Part~I).

\begin{lstlisting}
import tensorly as tl
from tensorly.decomposition import tucker
tl.set_backend("pytorch")

K, V = cache[layer_idx]                 # K: [B, H, S, D]
Kten = K[0].float()                     # [H, S, D] = heads x tokens x head_dim

r_h, r_s, r_d = 8, 128, 16              # per-mode ranks (<= H, S, D); r_h < 8 shares across heads
core, factors = tucker(Kten, rank=[r_h, r_s, r_d])
Khat = tl.tucker_to_tensor((core, factors))          # reconstruction, shape [H, S, D]
rel_err = (Khat - Kten).norm() / Kten.norm()

dense      = Kten.numel()
compressed = core.numel() + sum(f.numel() for f in factors)
print(f"Tucker rel_err={rel_err:.4f}  compression={dense / compressed:.2f}x")
\end{lstlisting}
Compare this compression-vs-error point against granularity~(1) (per-head SVD) at \emph{matched}
memory: if Tucker wins, the cross-head sharing is real structure worth exploiting.
(\code{tensor\_train} gives the TT variant; reconstruct with \code{tl.tt\_to\_tensor}.)

% ================================================================
\part*{Part VI \;---\; Interpreting your results}
\addcontentsline{toc}{section}{Part VI --- Interpreting your results}

\section{What to expect (and the honest lessons)}
\label{sec:expect}

\begin{enumerate}[leftmargin=1.5em,itemsep=2pt]
\item \term{Keys compress, values resist.} Expect key spectra to decay fast and value spectra to
be flatter; ShadowKV reports values are \emph{not} low-rank~\cite{shadowkv}, which is why it keeps
low-rank keys but offloads values. If you see this, it is a real, reportable finding --- and it
argues for compressing keys and values \emph{differently}.
\item \term{Depth and RoPE are open questions.} The literature disagrees on whether early or deep
layers are more compressible, and on whether RoPE raises or lowers key
rank~\cite{shadowkv,lorc}. Your clean single-model measurement is a genuine contribution --- just
state your setup (which model, context length, pre- vs post-RoPE).
\item \term{Naive low-rank can be \emph{slower}.} This is the most important lesson. If you store
factors and reconstruct dense $\hat K$ every step, you pay an SVD/matmul per step; on a small head
dim ($d=64$/$128$) the dense attention was already cheap, so reconstruction overhead can exceed
the savings, and TPOT \emph{goes up}. This is not a failure of your experiment --- it is precisely
why the field moves to (a) storing factors for \emph{memory} while accepting some latency, (b)
computing attention \emph{in the compressed domain} (no reconstruction), or (c) query-aware sparse
retrieval. Documenting the crossover ($r$, context length) at which low-rank helps or hurts is a
publishable result on its own.
\item \term{Reconstruction error underestimates task damage.} A rank that looks fine on
$\lVert K-\hat K\rVert$ and even on perplexity can still hurt GSM8K, because chained reasoning
compounds per-step error~\cite{gear}. Always keep a reasoning task in the loop.
\end{enumerate}

\begin{keybox}
These ``negative'' findings are the bridge to the research program: the memory-wall poster's
pipeline~\cite{poster} plus \emph{compressed-domain attention} (compute from factors, never
reconstruct) is exactly the fix for lesson~3. Your single-device evaluation is the evidence base
that motivates it.
\end{keybox}

% ================================================================
\part*{Part VII \;---\; Related work}
\addcontentsline{toc}{section}{Part VII --- Related work}

\section{The landscape (grouped, with code where public)}

\textbf{KV quantization:} KIVI~\cite{kivi} (2-bit; keys per-channel, values per-token;
\href{https://github.com/jy-yuan/KIVI}{code}); GEAR~\cite{gear} (quant + low-rank error + sparse
outliers; \href{https://github.com/opengear-project/GEAR}{code}); DecoQuant~\cite{decoquant}
(tensor-decomposition-assisted low-bit KV; \href{https://github.com/lpyhdzx/DecoQuant_code}{code}).
\\[2pt]
\textbf{Token eviction / attention sparsity:} H2O~\cite{h2o} (heavy hitters); SnapKV~\cite{snapkv};
Quest~\cite{quest} (query-aware page selection); RocketKV~\cite{rocketkv} (two-stage
eviction + sparse attention).
\\[2pt]
\textbf{Low-rank / tensor KV (closest to this project):} Palu~\cite{palu} (low-rank projection of
KV; \href{https://github.com/shadowpa0327/Palu}{code}); ShadowKV~\cite{shadowkv} (low-rank keys on
GPU, values offloaded; \href{https://github.com/ByteDance-Seed/ShadowKV}{code}); Eigen
Attention~\cite{eigenattn}; MLA / DeepSeek-V2~\cite{mla} (trained-in low-rank latent KV);
JoLT~\cite{jolt} (partial Tucker of the KV tensor + JL residual --- the explicit tensor version;
recent preprint).
\\[2pt]
\textbf{Attention-output-error theory (why not $\lVert K-\hat K\rVert$):} KDEformer~\cite{kdeformer};
Ada-KV~\cite{adakv} (error bounded by evicted attention mass); the total-variation / missed-mass
identity~\cite{tvattn} (recent preprint).
\\[2pt]
\textbf{Tensor foundations:} Robust PCA (low-rank + sparse)~\cite{rpca}; CP/Tucker~\cite{kolda};
Tensor-Train~\cite{tt}; low-rank + sparse weights (LoSparse)~\cite{losparse}; low-rank attention
(Linformer)~\cite{linformer}.
\\[2pt]
\textbf{Systems \& the base pipeline:} FlashAttention~\cite{flashattention} (the online-softmax
primitive); vLLM / PagedAttention~\cite{pagedattention} (the standard KV memory manager you are
implicitly compared against); the FSU pipelined-inference poster~\cite{poster}, HexGen~\cite{hexgen},
Jupiter~\cite{jupiter}.
\\[2pt]
\textbf{Evaluation:} WikiText-2~\cite{wikitext}; C4~\cite{c4}; LongBench~\cite{longbench};
RULER~\cite{ruler}; GSM8K~\cite{gsm8k}; the HF perplexity protocol~\cite{hfppl}.

\begin{tipbox}
Recent (2025--2026) preprints marked as such (JoLT, the TV/missed-mass identity) are verified to
exist on arXiv but are not yet peer-reviewed --- read the PDF before leaning on their exact
numbers. Everything else here is published.
\end{tipbox}

% ================================================================
\part*{Part VIII \;---\; Deliverables \& checklist}
\addcontentsline{toc}{section}{Part VIII --- Deliverables \& checklist}

\section{What to produce}
\begin{itemize}[leftmargin=1.5em,itemsep=1pt]
\item \textbf{Plots:} (1) singular-value / energy-vs-rank spectra, keys vs values, by layer;
(2) KL-divergence vs rank; (3) compression-vs-accuracy Pareto (per keys/values/both);
(4) TPOT and peak-memory vs rank and context length; (5) energy/token vs rank.
\item \textbf{Tables:} baseline metrics; the crossover $r$ where low-rank helps/hurts latency;
your method vs KIVI 2-bit vs full cache.
\item \textbf{One figure that tells the story:} accuracy retained (y) vs compression ratio (x),
with a second panel of latency vs compression --- this is the paper's money plot.
\end{itemize}

\section{Reproducibility checklist}
\begin{itemize}[leftmargin=1.5em,itemsep=0pt]
\item[$\square$] Fixed seeds; greedy decoding for all accuracy comparisons.
\item[$\square$] \code{cuda.synchronize} / CUDA events around every timing; warmup done.
\item[$\square$] Median and P95 reported over $\ge$50 iters.
\item[$\square$] KV bytes reported both analytically and measured.
\item[$\square$] Same tokenizer and context length across compared runs (perplexity is tokenizer-dependent).
\item[$\square$] Pre- vs post-RoPE stated for any spectrum plot.
\item[$\square$] A reasoning task (GSM8K) included, not only perplexity.
\item[$\square$] transformers / torch versions recorded in the writeup.
\end{itemize}

% ================================================================
\begin{thebibliography}{99}
\small
% base system
\bibitem{poster} H.~Zhou, X.~Zhang. \emph{Poster: Distributed Pipelined Collaborative Inference for LLMs on Edge Device Clusters.} ACM HotMobile 2026.
\bibitem{hexgen} Y.~Jiang et al. \emph{HexGen: Generative Inference of LLM over Heterogeneous Environment.} ICML 2024. arXiv:2311.11514.
\bibitem{jupiter} S.~Ye et al. \emph{Jupiter: Fast and Resource-Efficient Collaborative Inference of Generative LLMs on Edge Devices.} INFOCOM 2025. arXiv:2504.08242.
% KV quantization
\bibitem{kivi} Z.~Liu et al. \emph{KIVI: A Tuning-Free Asymmetric 2bit Quantization for KV Cache.} ICML 2024. arXiv:2402.02750.
\bibitem{gear} H.~Kang, Q.~Zhang, G.~Jeong, T.~Krishna, T.~Zhao. \emph{GEAR: An Efficient Error Reduction Framework for KV Cache Compression.} PMLR v262, 2024. arXiv:2403.05527.
\bibitem{decoquant} P.~Liu et al. \emph{Unlocking Data-free Low-bit Quantization with Matrix Decomposition for KV Cache Compression (DecoQuant).} ACL 2024. arXiv:2405.12591.
% eviction / sparsity
\bibitem{h2o} Z.~Zhang et al. \emph{H2O: Heavy-Hitter Oracle for Efficient Generative Inference of LLMs.} NeurIPS 2023. arXiv:2306.14048.
\bibitem{snapkv} Y.~Li et al. \emph{SnapKV: LLM Knows What You Are Looking For Before Generation.} arXiv:2404.14469, 2024.
\bibitem{quest} J.~Tang et al. \emph{Quest: Query-Aware Sparsity for Efficient Long-Context LLM Inference.} ICML 2024. arXiv:2406.10774.
\bibitem{rocketkv} P.~Behnam et al. \emph{RocketKV: Accelerating Long-Context LLM Inference via Two-Stage KV Cache Compression.} ICML 2025. arXiv:2502.14051.
% low-rank / tensor KV
\bibitem{palu} J.~Chang et al. \emph{Palu: Compressing KV-Cache with Low-Rank Projection.} ICLR 2025. arXiv:2407.21118.
\bibitem{shadowkv} H.~Sun et al. \emph{ShadowKV: KV Cache in Shadows for High-Throughput Long-Context LLM Inference.} ICML 2025. arXiv:2410.21465.
\bibitem{eigenattn} U.~Saxena et al. \emph{Eigen Attention: Attention in Low-Rank Space for KV Cache Compression.} EMNLP 2024 Findings. arXiv:2408.05646.
\bibitem{mla} DeepSeek-AI. \emph{DeepSeek-V2} (Multi-head Latent Attention: trained-in low-rank KV). arXiv:2405.04434, 2024.
\bibitem{jolt} \emph{A JoLT for the KV Cache: Near-Lossless KV Cache Compression via Joint Tucker and JL-Residual Allocation.} arXiv:2607.12550, 2026 (preprint).
% attention-output error theory
\bibitem{kdeformer} A.~Zandieh, I.~Han, M.~Daliri, A.~Karbasi. \emph{KDEformer: Accelerating Transformers via Kernel Density Estimation.} ICML 2023. arXiv:2302.02451.
\bibitem{adakv} Y.~Feng et al. \emph{Ada-KV: Optimizing KV Cache Eviction by Adaptive Budget Allocation.} arXiv:2407.11550, 2024.
\bibitem{tvattn} A.~Tzachristas et al. \emph{A Mathematical Theory of Top-k Sparse Attention via Total Variation Distance.} arXiv:2512.07647, 2025 (preprint).
% tensor foundations
\bibitem{rpca} E.~J.~Cand\`es, X.~Li, Y.~Ma, J.~Wright. \emph{Robust Principal Component Analysis?} J.~ACM 58(3):11, 2011.
\bibitem{kolda} T.~G.~Kolda, B.~W.~Bader. \emph{Tensor Decompositions and Applications.} SIAM Review 51(3):455--500, 2009.
\bibitem{tt} I.~V.~Oseledets. \emph{Tensor-Train Decomposition.} SIAM J.~Sci.~Comput. 33(5):2295--2317, 2011.
\bibitem{losparse} Y.~Li et al. \emph{LoSparse: Structured Compression of LLMs based on Low-Rank and Sparse Approximation.} ICML 2023. arXiv:2306.11222.
\bibitem{linformer} S.~Wang et al. \emph{Linformer: Self-Attention with Linear Complexity.} arXiv:2006.04768, 2020.
\bibitem{lorc} R.~Zhang et al. \emph{LoRC: Low-Rank Compression for LLMs KV Cache with a Progressive Compression Strategy.} arXiv:2410.03111, 2024.
% systems
\bibitem{flashattention} T.~Dao, D.~Y.~Fu, S.~Ermon, A.~Rudra, C.~R\'e. \emph{FlashAttention: Fast and Memory-Efficient Exact Attention with IO-Awareness.} NeurIPS 2022. arXiv:2205.14135.
\bibitem{pagedattention} W.~Kwon et al. \emph{Efficient Memory Management for LLM Serving with PagedAttention (vLLM).} SOSP 2023. arXiv:2309.06180.
% evaluation
\bibitem{wikitext} S.~Merity et al. \emph{Pointer Sentinel Mixture Models} (WikiText). arXiv:1609.07843, 2016.
\bibitem{c4} C.~Raffel et al. \emph{Exploring the Limits of Transfer Learning with a Unified Text-to-Text Transformer} (C4). arXiv:1910.10683, 2019.
\bibitem{longbench} Y.~Bai et al. \emph{LongBench: A Bilingual, Multitask Benchmark for Long Context Understanding.} ACL 2024. arXiv:2308.14508.
\bibitem{ruler} C.-P.~Hsieh et al. \emph{RULER: What's the Real Context Size of Your Long-Context Language Models?} COLM 2024. arXiv:2404.06654.
\bibitem{gsm8k} K.~Cobbe et al. \emph{Training Verifiers to Solve Math Word Problems (GSM8K).} arXiv:2110.14168, 2021.
\bibitem{hfppl} Hugging Face. \emph{Perplexity of fixed-length models.} \url{https://huggingface.co/docs/transformers/en/perplexity}.
\end{thebibliography}

\end{document}
